% Generated draft: requires mathematical and source review.
\documentclass[11pt]{article}
\usepackage[margin=25mm]{geometry}
\usepackage{fontspec}
\setmainfont{Latin Modern Roman}
\usepackage{amsmath,amssymb,mathtools,unicode-math}
\setmathfont{Latin Modern Math}
\usepackage{enumitem,needspace,xurl,hyperref}
\setlength{\parindent}{0pt}
\setlength{\parskip}{5pt}
\setlength{\emergencystretch}{3em}
\newcommand{\N}{\mathbb{N}}
\newcommand{\Z}{\mathbb{Z}}
\newcommand{\Q}{\mathbb{Q}}
\newcommand{\R}{\mathbb{R}}
\newcommand{\C}{\mathbb{C}}
\newcommand{\F}{\mathbb{F}}
\newcommand{\A}{\mathbb{A}}
\newcommand{\PP}{\mathbb{P}}
\newcommand{\E}{\mathbb{E}}
\newcommand{\eps}{\varepsilon}
\newcommand{\dd}{\,\mathrm{d}}
\begin{document}
\section*{UnsolvedMath ID: 30003936}
\textbf{Source problem code:} catalogue entry
\subsection{Definitions and mathematical statement}
The supplied registry entry preserves a modeling agenda rather than a recoverable standalone open problem. Its original questions concern a free-boundary problem with a random domain \(D(\omega)\) and random diffusion.

For interpreting the almost-sure qualification, let \((\Omega,\mathcal F,\mathbb P)\) denote an unspecified probability space, with \(\omega\in\Omega\) an outcome and \(D(\omega)\) the proposed domain associated with that outcome. These symbols describe the requested framework; the record does not supply a construction or a measurable space of admissible domains.

The preserved agenda asks for:
\begin{enumerate}
\item a model and hypotheses under which \(D(\omega)\) is well-defined almost surely, meaning that it is well-defined for every \(\omega\) in some \(\Omega_0\in\mathcal F\) with \(\mathbb P(\Omega_0)=1\);
\item an appropriate definition of the expectation of the random domain;
\item numerical methods for computing the solution and the domain's expectation and deviation.
\end{enumerate}
The meaning of ``well-defined,'' the governing equation, diffusion law, spatial setting, free-boundary conditions, solution concept, domain expectation, and deviation are unspecified. Consequently, these requests do not determine a precise existence, uniqueness, or numerical-convergence assertion.

The supplied statement-verification field says that the official contribution presents these questions as its agenda and then provides a well-posedness framework, expectation and deviation concepts, and a numerical method. It therefore classifies the extracted statement as \texttt{unrecoverable} and promotes no separate clean statement. This supplied verification has not been independently checked. Although the registry retains an \texttt{open} label, that label does not establish that the preserved agenda is unresolved.

The supplied dated literature assessment instead concerns the Dynamic Monge--Kantorovich flow: it describes existence beyond local time, global well-posedness, and convergence to the optimal-transport density and potential as open, and states that numerical evidence and Lyapunov candidates do not constitute a proof. The record gives no connection between that assessment and the random-domain agenda, so it cannot supply the missing mathematical formulation.
\subsection{Short English statement}
The retained questions ask how to model a random free boundary with random diffusion, define the domain's expectation and deviation, and compute these quantities numerically. The supplied verification identifies them as an agenda addressed within the source contribution, not a standalone unresolved problem. Essential model data are absent, and the accompanying Dynamic Monge--Kantorovich assessment concerns a different formulation. Human source review is required before a precise open problem can be stated.
\subsection{Sources}
\begin{itemize}
\item Local UnsolvedMath registry, permanent ID \(30003936\), source code \texttt{catalogue entry}, title ``Source-Resolved Random-Domain Modeling Agenda,'' in the Oberwolfach Reports Open Problems collection. The supplied record lists proposed year \(2018\) and names no proposer. Its original statement is retained for provenance; its clean-statement field is empty and its statement status is \texttt{unrecoverable}.
\item Supplied source citation: \emph{New Directions in Stochastic Optimisation} (2019), Oberwolfach Reports, pp.\ 2303--2384, DOI \texttt{10.4171/owr/2018/38}. \url{https://doi.org/10.4171/owr/2018/38}. The record attributes the questions to a workshop contribution but supplies no contribution title, author, or precise page within this range.
\item Supplied literature source: \emph{Towards a stationary Monge--Kantorovich dynamics}. \url{https://doi.org/10.1137/16M1098383}.
\item Supplied literature source: \emph{Oberwolfach Report 39/2018}. \url{https://doi.org/10.4171/owr/2018/39}. This differs from the report DOI in the primary source citation; its relevance to the random-domain questions is not established by the record.
\item Literature assessment date: \texttt{2026-08-21}. The supplied dated curation labels the Dynamic Monge--Kantorovich questions open and cautions that numerical evidence and Lyapunov candidates are not proofs. It explicitly requests independent verification of time-sensitive status and cited claims. This is supplied dated curation, not an independently verified literature assessment, and does not verify the status of the random-domain agenda.
\end{itemize}
\end{document}
